% Problem dossier: VI.02.P1
\subsection*{Problem VI.02.P1 --- Zariski topology versus product topology}
\label{prob:vi02-p1}

\paragraph{Problem.}
Assume \(k\) is infinite.  Describe the proper closed subsets of \(\mathbb A_k^1\) and use
this to describe the proper closed subsets of
\(\mathbb A_k^1\times\mathbb A_k^1\) equipped with the ordinary product topology.  Compare
that topology with the Zariski topology on \(\mathbb A_k^2\).  In particular, determine
whether the diagonal
\[
\Delta=\set{(t,t):t\in k}
\]
is closed in each topology.

\ifdefined\FullProblemDossiers
\paragraph{Why this problem matters.}
The underlying set \(k^2\) admits two natural-looking topologies, and they are not the same.
This is one of the first places where algebraic geometry defeats product-topology intuition.

\paragraph{Useful ideas before starting.}
On \(\mathbb A_k^1\), proper closed sets are finite.  In the product topology, start from a
point of the open complement of a proper closed set and place a cofinite product neighborhood
around it.

\paragraph{Guided hints.}
\textbf{Hint 1.} If \(F\) is proper product-closed, choose \((a,b)\notin F\) and a basic open
\(U\times V\subseteq F^c\) containing it.

\textbf{Hint 2.} Then \(F\subseteq (k\setminus U)\times k\;\cup\;k\times(k\setminus V)\), a
finite union of vertical and horizontal lines.

\textbf{Hint 3.} The diagonal is \(V(y-x)\) in the affine-plane Zariski topology.  To see it
is not product-closed, show every product neighborhood of an off-diagonal point meets it.

\paragraph{Solution.}
Because \(k\) is infinite, the proper closed subsets of \(\mathbb A_k^1\) are finite.  Let
\(F\) be a proper closed set in the product topology.  Choose \((a,b)\notin F\).  Since
\(F^c\) is open, there are cofinite sets \(U,V\subseteq k\) such that
\[
(a,b)\in U\times V\subseteq F^c.
\]
Hence
\[
F\subseteq (k\setminus U)\times k\;\cup\;k\times(k\setminus V),
\]
a finite union of vertical and horizontal lines.  Intersecting \(F\) with each such line,
whose subspace topology is again cofinite, shows that \(F\) is a finite union of whole
vertical lines, whole horizontal lines, and finitely many points.

Every basic product open \(U\times V\) is Zariski open in \(\mathbb A_k^2\): write
\(U=D(f(x))\), \(V=D(g(y))\), so
\[
U\times V=D(f(x)g(y)).
\]
Thus the product topology is coarser than the affine-plane Zariski topology.

The inclusion is strict.  The diagonal is Zariski closed because
\[
\Delta=V(y-x).
\]
Suppose it were product-closed.  Take \((a,b)\notin\Delta\).  Any product neighborhood
\(U\times V\) of \((a,b)\) has \(U,V\) cofinite and therefore \(U\cap V\neq\varnothing\).
For \(t\in U\cap V\), the diagonal point \((t,t)\) lies in \(U\times V\).  Hence the
complement of \(\Delta\) is not product-open, so \(\Delta\) is not product-closed.

\paragraph{What happened mathematically.}
The affine-plane Zariski topology sees equations relating different coordinates, such as
\(y=x\).  The product topology only knows the two one-dimensional topologies independently.

\paragraph{Extensions.}
\begin{enumerate}
\item Prove the same strictness for \(\mathbb A_k^n\) versus the \(n\)-fold product topology
for every \(n\ge2\).
\item Classify the closure of the diagonal in the product topology.
\item Determine what changes when \(k\) is finite.
\end{enumerate}

\paragraph{Provenance.}
Recovered and normalized from legacy DG4 affine Problem 1.  The infinite-field hypothesis is made explicit because the legacy solution uses
the cofinite description of \(\mathbb A^1_k\).
\fi
